%% Copyright (C) 2026 Barak A. Pearlmutter

\chapter{Threads} \label{threads}

Oaklisp has two kinds of concurrency.  \emph{Processes} are
lightweight: any number of them share one or more operating-system
threads under a scheduler of Oaklisp's own, switching at
\df{pause} and when preempted, and they are what \df{future},
\df{delay} and the semaphores of this chapter are built on.
\emph{Threads} are operating-system threads: with an emulator built
with thread support (\texttt{configure --enable-threads}) Oaklisp
code runs on several of them at once, and a program may make one of
its own with \df{make-thread}.

Each thread runs one process from a thunk to its end.  That process
starts with a copy of its creator's fluid bindings and its creator's
winding stack, so \df{bind} and \df{wind-protect} nest as they would
in a call; global variables and all other objects are shared.  The
runtime's own tables (the symbol table, the locales, the streams) are
safe to use from several threads at once, and \df{format} writes a
call's whole output to a stream at once; a hash table of the program's
own is safe to read while another thread writes it, but two writers
need a mutex.

\section{Threads}

\op{make-thread}{thunk $[$name$]$}
\doc{Starts an operating-system thread running \emph{thunk} and
returns its \df{thread} object at once.  \emph{Name} is for
messages and \df{thread-name}; by default the threads are numbered.
An error if the emulator has no thread support.}

\op{thread-join}{thread $[$milliseconds$]$}
\doc{Waits for \emph{thread} to finish and returns the value of its
thunk; joined again, returns it again.  If the thunk signaled an
error that nothing in the thread handled, \df{thread-join} signals an
error in the joiner, carrying the thread's.  With \emph{milliseconds},
waits at most that long and returns \texttt{\#f} if the thread is
still running.  While it waits, the joiner's own operating-system
thread is parked: no other process runs on it meanwhile.}

\op{thread-alive?}{thread}
\doc{Whether \emph{thread} has not finished yet.}

\op{thread-name}{thread}
\doc{The name \emph{thread} was made with.}

\op{current-thread}{}
\doc{The \df{thread} object of the thread this is called on; the
main thread's is named \texttt{main}, a scheduler thread's
\texttt{scheduler}.}

\op{thread-count}{}
\doc{The number of operating-system threads running Oaklisp code.}

\op{thread-sleep}{milliseconds}
\doc{Suspends the calling thread for \emph{milliseconds}.}

\section{Mutexes}

A mutex may be held by one process at a time.  It is held by a
\emph{process}: a process that is preempted inside a critical section
and resumed on another thread still holds it, and only the holder may
release it.  A mutex is reentrant: the holder may acquire it again,
and must release it as many times.

\op{make-named-mutex}{name}
\doc{A new mutex, named for messages; \texttt{(make mutex)} makes an
anonymous one.}

\op{acquire-mutex}{mutex}
\doc{Waits until \emph{mutex} is free and takes it.}

\op{release-mutex}{mutex}
\doc{Releases \emph{mutex}; an error if the caller does not hold it.}

\mc{with-mutex}{mutex \dt body}
\doc{Evaluates \emph{body} holding \emph{mutex}, releasing it however
the body is left, including by a throw.}

\section{Condition Variables}

\op{condition-variable}{}
\doc{The type; \texttt{(make condition-variable)} makes one.  Used
with a mutex that protects the state being waited for.}

\op{condition-wait}{condition-variable mutex $[$milliseconds$]$}
\doc{The caller must hold \emph{mutex}.  Releases it, waits until the
condition variable is notified (or for at most \emph{milliseconds}),
takes \emph{mutex} again and returns \texttt{\#t}, or \texttt{\#f}
after a timeout.  Wake-ups can be spurious: test the condition in a
loop around the wait.}

\op{condition-notify}{condition-variable}
\doc{Wakes one waiter, if any; whether there was one.}

\op{condition-notify-all}{condition-variable}
\doc{Wakes every waiter; how many.}

\section{Atomic Operations}

These act on a cell given as a locative, as \texttt{(make-locative
(car x))} or \texttt{(make-locative variable)}, and compare with
\df{eq?}, so they are for fixnums and objects, not bignums.

\op{compare-and-swap}{locative old new}
\doc{If the cell holds \emph{old} (by \df{eq?}), stores \emph{new}
and returns \emph{old}; otherwise returns what the cell holds,
unchanged.  Atomic with respect to every other thread.}

\op{fetch-and-add}{locative n}
\doc{Adds \emph{n} to the fixnum in the cell and returns the value it
had.}

\section{Examples}

Two programs in \texttt{doc/examples}: \texttt{parallel-map.oak}, a
\df{map} that runs each application on a thread of its own, and
\texttt{producer-consumer.oak}, a bounded queue with a mutex and two
condition variables between a producer thread and a consumer.
